% ============================================================================
%  preambule.tex : paquets, mise en page, notations et environnements
%  (inclus dans main.tex par % ============================================================================
%  preambule.tex : paquets, mise en page, notations et environnements
%  (inclus dans main.tex par % ============================================================================
%  preambule.tex : paquets, mise en page, notations et environnements
%  (inclus dans main.tex par % ============================================================================
%  preambule.tex : paquets, mise en page, notations et environnements
%  (inclus dans main.tex par \input{preambule})
% ============================================================================

% ----------------------------------------------------------------------------
%  Polices et langue
% ----------------------------------------------------------------------------
\usepackage[T1]{fontenc}
\usepackage[utf8]{inputenc}     % à SUPPRIMER si lualatex/xelatex
\usepackage[french]{babel}
\usepackage{lmodern}
\usepackage{microtype}          % justification nettement meilleure

% ----------------------------------------------------------------------------
%  Géométrie
% ----------------------------------------------------------------------------
% 11 pt et 3 cm de marges latérales : environ 80 caractères par ligne,
% texte dense mais lisible.
\usepackage[a4paper,
            left=2.0cm, right=2.0cm, top=2.5cm, bottom=2.5cm,
            headheight=15pt]{geometry}

% ----------------------------------------------------------------------------
%  Mathématiques
% ----------------------------------------------------------------------------
\usepackage{amsmath, amssymb, amsthm, mathtools}
\usepackage{bm}
\usepackage{mathrsfs}

% ----------------------------------------------------------------------------
%  Flottants, tableaux, listes
% ----------------------------------------------------------------------------
\usepackage{graphicx}
\usepackage{float}
\usepackage{booktabs}
\usepackage{caption}
\usepackage{subcaption}         % APRÈS caption
\usepackage[inline]{enumitem}

% ----------------------------------------------------------------------------
%  Dessins
% ----------------------------------------------------------------------------
\usepackage{tikz}
\usetikzlibrary{arrows.meta, positioning, calc, angles, quotes,
    decorations.pathmorphing, patterns}

% ----------------------------------------------------------------------------
%  Couleurs, encadrés et liens
% ----------------------------------------------------------------------------
\usepackage{xcolor}
\definecolor{accent}{HTML}{1F3A5F}
\usepackage[most]{tcolorbox}

% Ligne « institution » de la page de titre.
\newcommand{\kicker}[1]{{\scshape\large\textcolor{accent}{#1}}}

\usepackage[colorlinks=true, linkcolor=accent, citecolor=accent,
            urlcolor=accent]{hyperref}
\usepackage{cleveref}           % IMPÉRATIVEMENT après hyperref

% ============================================================================
%  MISE EN PAGE
% ============================================================================

% --- 1. Titres ---------------------------------------------------------------
\usepackage{titlesec}

% Chapitre : « Chapitre » en petites capitales, gros numéro, titre
% en dessous, filet pour fermer.
\titleformat{\chapter}[display]
    {\normalfont}
    {%
        {\color{accent!55}\large\scshape\chaptertitlename}\;\;
        {\color{accent}\fontsize{52}{52}\selectfont\bfseries\thechapter}%
    }
    {10pt}
    {\color{accent}\Huge\bfseries}
    [\vspace{4mm}{\color{accent!40}\titlerule[1.4pt]}]
\titlespacing*{\chapter}{0pt}{20pt}{34pt}

% Structure : Chapitre -> I. -> 1. -> 1.1 -> a)
\setcounter{secnumdepth}{4}
\setcounter{tocdepth}{3}
\renewcommand{\thesection}{\Roman{section}.}
\renewcommand{\thesubsection}{\arabic{subsection}.}
\renewcommand{\thesubsubsection}{%
    \arabic{subsection}.\arabic{subsubsection}}
\renewcommand{\theparagraph}{\alph{paragraph})}

\titlespacing*{\section}{0pt}{2.2ex plus 1ex minus .2ex}{1.4ex plus .2ex}
\titlespacing*{\subsection}{0pt}{1.8ex plus 1ex minus .2ex}{1.0ex plus .2ex}
% a) : titre sur sa propre ligne
\titleformat{\paragraph}
    {\normalfont\normalsize\bfseries}{\theparagraph}{0.5em}{}
\titlespacing*{\paragraph}{0pt}{1.4ex plus .5ex minus .2ex}{0.6ex}

% --- 2. Flottants ------------------------------------------------------------
\renewcommand{\topfraction}{0.90}
\renewcommand{\bottomfraction}{0.75}
\renewcommand{\textfraction}{0.10}
\renewcommand{\floatpagefraction}{0.75}
\setcounter{topnumber}{3}
\setcounter{bottomnumber}{2}
\setcounter{totalnumber}{5}

% --- 3. Coupures de page -----------------------------------------------------
\widowpenalty=10000
\clubpenalty=10000
\raggedbottom

% --- 4. Listes et légendes ---------------------------------------------------
\setlist{itemsep=2pt, parsep=0pt, topsep=4pt}
\captionsetup{font=small, labelfont={bf,color=accent}, skip=6pt}

% --- 5. Accroche : phrase clé mise en exergue --------------------------------
% Usage : \begin{accroche} ... \end{accroche}
\newenvironment{accroche}[1][accent]{%
    \par\medskip\noindent
    \begingroup
    \color{#1}\itshape\large
    \rightskip=1.5em \leftskip=1.5em
    \parindent=0pt
    \hrule height 0.9pt\vspace{7pt}%
}{%
    \par\vspace{7pt}\hrule height 0.9pt
    \endgroup
    \par\medskip
}

% ============================================================================
%  NOTATIONS
% ============================================================================
% Vecteur : minuscule ou grecque, en gras italique
% (ex. \vect{u}, \vect{\omega}).
\newcommand{\vect}[1]{\bm{#1}}
% Matrice / tenseur : majuscule en gras droit (ex. \mat{R}, \mat{I}).
\newcommand{\mat}[1]{\mathbf{#1}}
% Quaternion : gras, souligné (ex. \quat{q} = (q_0, \vect{q})).
\newcommand{\quat}[1]{\underline{\bm{#1}}}
% Produit quaternionique : étoile centrée (\otimes est réservé au
% produit tensoriel).
\newcommand{\quatmul}{\ast}
% Repères : \rep{s} (sol), \rep{b} (solide).
\newcommand{\rep}[1]{\mathcal{R}_{#1}}
% Matrice antisymétrique associée à u : [u]_x v = u \wedge v
\newcommand{\hatx}[1]{[#1]_{\times}}
% Ensembles et opérateurs
\newcommand{\R}{\mathbb{R}}
\newcommand{\SO}{\mathrm{SO}(3)}
\newcommand{\Id}{\mat{I}}
\newcommand{\dd}{\mathrm{d}}
\newcommand{\ddt}{\frac{\dd}{\dd t}}
\newcommand{\pd}[2]{\frac{\partial #1}{\partial #2}}
\DeclareMathOperator{\rang}{rang}
\DeclareMathOperator{\Ker}{Ker}
\DeclareMathOperator{\Ima}{Im}

% ============================================================================
%  ENVIRONNEMENTS (numérotés par chapitre)
% ============================================================================
\theoremstyle{plain}
\newtheorem{theoreme}{Théorème}[chapter]
\newtheorem{corollaire}[theoreme]{Corollaire}
\newtheorem{proposition}[theoreme]{Proposition}
\newtheorem{lemme}[theoreme]{Lemme}
\theoremstyle{definition}
\newtheorem{definition}[theoreme]{Définition}
\newtheorem{hypotheses}[theoreme]{Hypothèses}
\newtheorem{exemple}[theoreme]{Exemple}
\theoremstyle{remark}
\newtheorem{remarque}[theoreme]{Remarque}
\newtheorem{note}[theoreme]{Note}
\renewcommand{\proofname}{Démonstration}
% Démonstration détaillée et preuve courte : même mise en forme
% (titre en italique, carré en fin).
\newenvironment{demonstration}[1][Démonstration]
    {\begin{proof}[#1]}{\end{proof}}
\newenvironment{preuve}[1][Preuve]
    {\begin{proof}[#1]}{\end{proof}}

% Couleurs des cadres : théorème en rouge, définition et hypothèses
% en vert, proposition, corollaire et lemme en bleu.
\definecolor{coulthm}{HTML}{A8322D}
\definecolor{couldef}{HTML}{2E7D32}
\definecolor{coulprop}{HTML}{1F5A9E}
\tcbset{
    cadre/.style={enhanced,colback=#1!7,colframe=#1!80,
        boxrule=0.6pt,leftrule=3pt,arc=1pt,left=6pt,right=6pt,
        before skip=6pt,after skip=6pt}}
\tcolorboxenvironment{theoreme}{cadre=coulthm}
\tcolorboxenvironment{definition}{cadre=couldef}
\tcolorboxenvironment{hypotheses}{cadre=couldef}
\tcolorboxenvironment{proposition}{cadre=coulprop}
\tcolorboxenvironment{corollaire}{cadre=coulprop}
\tcolorboxenvironment{lemme}{cadre=coulprop}
% Note : notation ou commentaire personnel, cadre gris.
\tcolorboxenvironment{note}{cadre=gray}

% Renvois \cref en français
\crefname{theoreme}{théorème}{théorèmes}
\Crefname{theoreme}{Théorème}{Théorèmes}
\crefname{corollaire}{corollaire}{corollaires}
\Crefname{corollaire}{Corollaire}{Corollaires}
\crefname{proposition}{proposition}{propositions}
\Crefname{proposition}{Proposition}{Propositions}
\crefname{lemme}{lemme}{lemmes}
\Crefname{lemme}{Lemme}{Lemmes}
\crefname{definition}{définition}{définitions}
\Crefname{definition}{Définition}{Définitions}
\crefname{hypotheses}{hypothèses}{hypothèses}
\Crefname{hypotheses}{Hypothèses}{Hypothèses}
\crefname{exemple}{exemple}{exemples}
\Crefname{exemple}{Exemple}{Exemples}
\crefname{remarque}{remarque}{remarques}
\Crefname{remarque}{Remarque}{Remarques}
\crefname{note}{note}{notes}
\Crefname{note}{Note}{Notes}
\newcommand{\creflastconjunction}{ et~}
)
% ============================================================================

% ----------------------------------------------------------------------------
%  Polices et langue
% ----------------------------------------------------------------------------
\usepackage[T1]{fontenc}
\usepackage[utf8]{inputenc}     % à SUPPRIMER si lualatex/xelatex
\usepackage[french]{babel}
\usepackage{lmodern}
\usepackage{microtype}          % justification nettement meilleure

% ----------------------------------------------------------------------------
%  Géométrie
% ----------------------------------------------------------------------------
% 11 pt et 3 cm de marges latérales : environ 80 caractères par ligne,
% texte dense mais lisible.
\usepackage[a4paper,
            left=2.0cm, right=2.0cm, top=2.5cm, bottom=2.5cm,
            headheight=15pt]{geometry}

% ----------------------------------------------------------------------------
%  Mathématiques
% ----------------------------------------------------------------------------
\usepackage{amsmath, amssymb, amsthm, mathtools}
\usepackage{bm}
\usepackage{mathrsfs}

% ----------------------------------------------------------------------------
%  Flottants, tableaux, listes
% ----------------------------------------------------------------------------
\usepackage{graphicx}
\usepackage{float}
\usepackage{booktabs}
\usepackage{caption}
\usepackage{subcaption}         % APRÈS caption
\usepackage[inline]{enumitem}

% ----------------------------------------------------------------------------
%  Dessins
% ----------------------------------------------------------------------------
\usepackage{tikz}
\usetikzlibrary{arrows.meta, positioning, calc, angles, quotes,
    decorations.pathmorphing, patterns}

% ----------------------------------------------------------------------------
%  Couleurs, encadrés et liens
% ----------------------------------------------------------------------------
\usepackage{xcolor}
\definecolor{accent}{HTML}{1F3A5F}
\usepackage[most]{tcolorbox}

% Ligne « institution » de la page de titre.
\newcommand{\kicker}[1]{{\scshape\large\textcolor{accent}{#1}}}

\usepackage[colorlinks=true, linkcolor=accent, citecolor=accent,
            urlcolor=accent]{hyperref}
\usepackage{cleveref}           % IMPÉRATIVEMENT après hyperref

% ============================================================================
%  MISE EN PAGE
% ============================================================================

% --- 1. Titres ---------------------------------------------------------------
\usepackage{titlesec}

% Chapitre : « Chapitre » en petites capitales, gros numéro, titre
% en dessous, filet pour fermer.
\titleformat{\chapter}[display]
    {\normalfont}
    {%
        {\color{accent!55}\large\scshape\chaptertitlename}\;\;
        {\color{accent}\fontsize{52}{52}\selectfont\bfseries\thechapter}%
    }
    {10pt}
    {\color{accent}\Huge\bfseries}
    [\vspace{4mm}{\color{accent!40}\titlerule[1.4pt]}]
\titlespacing*{\chapter}{0pt}{20pt}{34pt}

% Structure : Chapitre -> I. -> 1. -> 1.1 -> a)
\setcounter{secnumdepth}{4}
\setcounter{tocdepth}{3}
\renewcommand{\thesection}{\Roman{section}.}
\renewcommand{\thesubsection}{\arabic{subsection}.}
\renewcommand{\thesubsubsection}{%
    \arabic{subsection}.\arabic{subsubsection}}
\renewcommand{\theparagraph}{\alph{paragraph})}

\titlespacing*{\section}{0pt}{2.2ex plus 1ex minus .2ex}{1.4ex plus .2ex}
\titlespacing*{\subsection}{0pt}{1.8ex plus 1ex minus .2ex}{1.0ex plus .2ex}
% a) : titre sur sa propre ligne
\titleformat{\paragraph}
    {\normalfont\normalsize\bfseries}{\theparagraph}{0.5em}{}
\titlespacing*{\paragraph}{0pt}{1.4ex plus .5ex minus .2ex}{0.6ex}

% --- 2. Flottants ------------------------------------------------------------
\renewcommand{\topfraction}{0.90}
\renewcommand{\bottomfraction}{0.75}
\renewcommand{\textfraction}{0.10}
\renewcommand{\floatpagefraction}{0.75}
\setcounter{topnumber}{3}
\setcounter{bottomnumber}{2}
\setcounter{totalnumber}{5}

% --- 3. Coupures de page -----------------------------------------------------
\widowpenalty=10000
\clubpenalty=10000
\raggedbottom

% --- 4. Listes et légendes ---------------------------------------------------
\setlist{itemsep=2pt, parsep=0pt, topsep=4pt}
\captionsetup{font=small, labelfont={bf,color=accent}, skip=6pt}

% --- 5. Accroche : phrase clé mise en exergue --------------------------------
% Usage : \begin{accroche} ... \end{accroche}
\newenvironment{accroche}[1][accent]{%
    \par\medskip\noindent
    \begingroup
    \color{#1}\itshape\large
    \rightskip=1.5em \leftskip=1.5em
    \parindent=0pt
    \hrule height 0.9pt\vspace{7pt}%
}{%
    \par\vspace{7pt}\hrule height 0.9pt
    \endgroup
    \par\medskip
}

% ============================================================================
%  NOTATIONS
% ============================================================================
% Vecteur : minuscule ou grecque, en gras italique
% (ex. \vect{u}, \vect{\omega}).
\newcommand{\vect}[1]{\bm{#1}}
% Matrice / tenseur : majuscule en gras droit (ex. \mat{R}, \mat{I}).
\newcommand{\mat}[1]{\mathbf{#1}}
% Quaternion : gras, souligné (ex. \quat{q} = (q_0, \vect{q})).
\newcommand{\quat}[1]{\underline{\bm{#1}}}
% Produit quaternionique : étoile centrée (\otimes est réservé au
% produit tensoriel).
\newcommand{\quatmul}{\ast}
% Repères : \rep{s} (sol), \rep{b} (solide).
\newcommand{\rep}[1]{\mathcal{R}_{#1}}
% Matrice antisymétrique associée à u : [u]_x v = u \wedge v
\newcommand{\hatx}[1]{[#1]_{\times}}
% Ensembles et opérateurs
\newcommand{\R}{\mathbb{R}}
\newcommand{\SO}{\mathrm{SO}(3)}
\newcommand{\Id}{\mat{I}}
\newcommand{\dd}{\mathrm{d}}
\newcommand{\ddt}{\frac{\dd}{\dd t}}
\newcommand{\pd}[2]{\frac{\partial #1}{\partial #2}}
\DeclareMathOperator{\rang}{rang}
\DeclareMathOperator{\Ker}{Ker}
\DeclareMathOperator{\Ima}{Im}

% ============================================================================
%  ENVIRONNEMENTS (numérotés par chapitre)
% ============================================================================
\theoremstyle{plain}
\newtheorem{theoreme}{Théorème}[chapter]
\newtheorem{corollaire}[theoreme]{Corollaire}
\newtheorem{proposition}[theoreme]{Proposition}
\newtheorem{lemme}[theoreme]{Lemme}
\theoremstyle{definition}
\newtheorem{definition}[theoreme]{Définition}
\newtheorem{hypotheses}[theoreme]{Hypothèses}
\newtheorem{exemple}[theoreme]{Exemple}
\theoremstyle{remark}
\newtheorem{remarque}[theoreme]{Remarque}
\newtheorem{note}[theoreme]{Note}
\renewcommand{\proofname}{Démonstration}
% Démonstration détaillée et preuve courte : même mise en forme
% (titre en italique, carré en fin).
\newenvironment{demonstration}[1][Démonstration]
    {\begin{proof}[#1]}{\end{proof}}
\newenvironment{preuve}[1][Preuve]
    {\begin{proof}[#1]}{\end{proof}}

% Couleurs des cadres : théorème en rouge, définition et hypothèses
% en vert, proposition, corollaire et lemme en bleu.
\definecolor{coulthm}{HTML}{A8322D}
\definecolor{couldef}{HTML}{2E7D32}
\definecolor{coulprop}{HTML}{1F5A9E}
\tcbset{
    cadre/.style={enhanced,colback=#1!7,colframe=#1!80,
        boxrule=0.6pt,leftrule=3pt,arc=1pt,left=6pt,right=6pt,
        before skip=6pt,after skip=6pt}}
\tcolorboxenvironment{theoreme}{cadre=coulthm}
\tcolorboxenvironment{definition}{cadre=couldef}
\tcolorboxenvironment{hypotheses}{cadre=couldef}
\tcolorboxenvironment{proposition}{cadre=coulprop}
\tcolorboxenvironment{corollaire}{cadre=coulprop}
\tcolorboxenvironment{lemme}{cadre=coulprop}
% Note : notation ou commentaire personnel, cadre gris.
\tcolorboxenvironment{note}{cadre=gray}

% Renvois \cref en français
\crefname{theoreme}{théorème}{théorèmes}
\Crefname{theoreme}{Théorème}{Théorèmes}
\crefname{corollaire}{corollaire}{corollaires}
\Crefname{corollaire}{Corollaire}{Corollaires}
\crefname{proposition}{proposition}{propositions}
\Crefname{proposition}{Proposition}{Propositions}
\crefname{lemme}{lemme}{lemmes}
\Crefname{lemme}{Lemme}{Lemmes}
\crefname{definition}{définition}{définitions}
\Crefname{definition}{Définition}{Définitions}
\crefname{hypotheses}{hypothèses}{hypothèses}
\Crefname{hypotheses}{Hypothèses}{Hypothèses}
\crefname{exemple}{exemple}{exemples}
\Crefname{exemple}{Exemple}{Exemples}
\crefname{remarque}{remarque}{remarques}
\Crefname{remarque}{Remarque}{Remarques}
\crefname{note}{note}{notes}
\Crefname{note}{Note}{Notes}
\newcommand{\creflastconjunction}{ et~}
)
% ============================================================================

% ----------------------------------------------------------------------------
%  Polices et langue
% ----------------------------------------------------------------------------
\usepackage[T1]{fontenc}
\usepackage[utf8]{inputenc}     % à SUPPRIMER si lualatex/xelatex
\usepackage[french]{babel}
\usepackage{lmodern}
\usepackage{microtype}          % justification nettement meilleure

% ----------------------------------------------------------------------------
%  Géométrie
% ----------------------------------------------------------------------------
% 11 pt et 3 cm de marges latérales : environ 80 caractères par ligne,
% texte dense mais lisible.
\usepackage[a4paper,
            left=2.0cm, right=2.0cm, top=2.5cm, bottom=2.5cm,
            headheight=15pt]{geometry}

% ----------------------------------------------------------------------------
%  Mathématiques
% ----------------------------------------------------------------------------
\usepackage{amsmath, amssymb, amsthm, mathtools}
\usepackage{bm}
\usepackage{mathrsfs}

% ----------------------------------------------------------------------------
%  Flottants, tableaux, listes
% ----------------------------------------------------------------------------
\usepackage{graphicx}
\usepackage{float}
\usepackage{booktabs}
\usepackage{caption}
\usepackage{subcaption}         % APRÈS caption
\usepackage[inline]{enumitem}

% ----------------------------------------------------------------------------
%  Dessins
% ----------------------------------------------------------------------------
\usepackage{tikz}
\usetikzlibrary{arrows.meta, positioning, calc, angles, quotes,
    decorations.pathmorphing, patterns}

% ----------------------------------------------------------------------------
%  Couleurs, encadrés et liens
% ----------------------------------------------------------------------------
\usepackage{xcolor}
\definecolor{accent}{HTML}{1F3A5F}
\usepackage[most]{tcolorbox}

% Ligne « institution » de la page de titre.
\newcommand{\kicker}[1]{{\scshape\large\textcolor{accent}{#1}}}

\usepackage[colorlinks=true, linkcolor=accent, citecolor=accent,
            urlcolor=accent]{hyperref}
\usepackage{cleveref}           % IMPÉRATIVEMENT après hyperref

% ============================================================================
%  MISE EN PAGE
% ============================================================================

% --- 1. Titres ---------------------------------------------------------------
\usepackage{titlesec}

% Chapitre : « Chapitre » en petites capitales, gros numéro, titre
% en dessous, filet pour fermer.
\titleformat{\chapter}[display]
    {\normalfont}
    {%
        {\color{accent!55}\large\scshape\chaptertitlename}\;\;
        {\color{accent}\fontsize{52}{52}\selectfont\bfseries\thechapter}%
    }
    {10pt}
    {\color{accent}\Huge\bfseries}
    [\vspace{4mm}{\color{accent!40}\titlerule[1.4pt]}]
\titlespacing*{\chapter}{0pt}{20pt}{34pt}

% Structure : Chapitre -> I. -> 1. -> 1.1 -> a)
\setcounter{secnumdepth}{4}
\setcounter{tocdepth}{3}
\renewcommand{\thesection}{\Roman{section}.}
\renewcommand{\thesubsection}{\arabic{subsection}.}
\renewcommand{\thesubsubsection}{%
    \arabic{subsection}.\arabic{subsubsection}}
\renewcommand{\theparagraph}{\alph{paragraph})}

\titlespacing*{\section}{0pt}{2.2ex plus 1ex minus .2ex}{1.4ex plus .2ex}
\titlespacing*{\subsection}{0pt}{1.8ex plus 1ex minus .2ex}{1.0ex plus .2ex}
% a) : titre sur sa propre ligne
\titleformat{\paragraph}
    {\normalfont\normalsize\bfseries}{\theparagraph}{0.5em}{}
\titlespacing*{\paragraph}{0pt}{1.4ex plus .5ex minus .2ex}{0.6ex}

% --- 2. Flottants ------------------------------------------------------------
\renewcommand{\topfraction}{0.90}
\renewcommand{\bottomfraction}{0.75}
\renewcommand{\textfraction}{0.10}
\renewcommand{\floatpagefraction}{0.75}
\setcounter{topnumber}{3}
\setcounter{bottomnumber}{2}
\setcounter{totalnumber}{5}

% --- 3. Coupures de page -----------------------------------------------------
\widowpenalty=10000
\clubpenalty=10000
\raggedbottom

% --- 4. Listes et légendes ---------------------------------------------------
\setlist{itemsep=2pt, parsep=0pt, topsep=4pt}
\captionsetup{font=small, labelfont={bf,color=accent}, skip=6pt}

% --- 5. Accroche : phrase clé mise en exergue --------------------------------
% Usage : \begin{accroche} ... \end{accroche}
\newenvironment{accroche}[1][accent]{%
    \par\medskip\noindent
    \begingroup
    \color{#1}\itshape\large
    \rightskip=1.5em \leftskip=1.5em
    \parindent=0pt
    \hrule height 0.9pt\vspace{7pt}%
}{%
    \par\vspace{7pt}\hrule height 0.9pt
    \endgroup
    \par\medskip
}

% ============================================================================
%  NOTATIONS
% ============================================================================
% Vecteur : minuscule ou grecque, en gras italique
% (ex. \vect{u}, \vect{\omega}).
\newcommand{\vect}[1]{\bm{#1}}
% Matrice / tenseur : majuscule en gras droit (ex. \mat{R}, \mat{I}).
\newcommand{\mat}[1]{\mathbf{#1}}
% Quaternion : gras, souligné (ex. \quat{q} = (q_0, \vect{q})).
\newcommand{\quat}[1]{\underline{\bm{#1}}}
% Produit quaternionique : étoile centrée (\otimes est réservé au
% produit tensoriel).
\newcommand{\quatmul}{\ast}
% Repères : \rep{s} (sol), \rep{b} (solide).
\newcommand{\rep}[1]{\mathcal{R}_{#1}}
% Matrice antisymétrique associée à u : [u]_x v = u \wedge v
\newcommand{\hatx}[1]{[#1]_{\times}}
% Ensembles et opérateurs
\newcommand{\R}{\mathbb{R}}
\newcommand{\SO}{\mathrm{SO}(3)}
\newcommand{\Id}{\mat{I}}
\newcommand{\dd}{\mathrm{d}}
\newcommand{\ddt}{\frac{\dd}{\dd t}}
\newcommand{\pd}[2]{\frac{\partial #1}{\partial #2}}
\DeclareMathOperator{\rang}{rang}
\DeclareMathOperator{\Ker}{Ker}
\DeclareMathOperator{\Ima}{Im}

% ============================================================================
%  ENVIRONNEMENTS (numérotés par chapitre)
% ============================================================================
\theoremstyle{plain}
\newtheorem{theoreme}{Théorème}[chapter]
\newtheorem{corollaire}[theoreme]{Corollaire}
\newtheorem{proposition}[theoreme]{Proposition}
\newtheorem{lemme}[theoreme]{Lemme}
\theoremstyle{definition}
\newtheorem{definition}[theoreme]{Définition}
\newtheorem{hypotheses}[theoreme]{Hypothèses}
\newtheorem{exemple}[theoreme]{Exemple}
\theoremstyle{remark}
\newtheorem{remarque}[theoreme]{Remarque}
\newtheorem{note}[theoreme]{Note}
\renewcommand{\proofname}{Démonstration}
% Démonstration détaillée et preuve courte : même mise en forme
% (titre en italique, carré en fin).
\newenvironment{demonstration}[1][Démonstration]
    {\begin{proof}[#1]}{\end{proof}}
\newenvironment{preuve}[1][Preuve]
    {\begin{proof}[#1]}{\end{proof}}

% Couleurs des cadres : théorème en rouge, définition et hypothèses
% en vert, proposition, corollaire et lemme en bleu.
\definecolor{coulthm}{HTML}{A8322D}
\definecolor{couldef}{HTML}{2E7D32}
\definecolor{coulprop}{HTML}{1F5A9E}
\tcbset{
    cadre/.style={enhanced,colback=#1!7,colframe=#1!80,
        boxrule=0.6pt,leftrule=3pt,arc=1pt,left=6pt,right=6pt,
        before skip=6pt,after skip=6pt}}
\tcolorboxenvironment{theoreme}{cadre=coulthm}
\tcolorboxenvironment{definition}{cadre=couldef}
\tcolorboxenvironment{hypotheses}{cadre=couldef}
\tcolorboxenvironment{proposition}{cadre=coulprop}
\tcolorboxenvironment{corollaire}{cadre=coulprop}
\tcolorboxenvironment{lemme}{cadre=coulprop}
% Note : notation ou commentaire personnel, cadre gris.
\tcolorboxenvironment{note}{cadre=gray}

% Renvois \cref en français
\crefname{theoreme}{théorème}{théorèmes}
\Crefname{theoreme}{Théorème}{Théorèmes}
\crefname{corollaire}{corollaire}{corollaires}
\Crefname{corollaire}{Corollaire}{Corollaires}
\crefname{proposition}{proposition}{propositions}
\Crefname{proposition}{Proposition}{Propositions}
\crefname{lemme}{lemme}{lemmes}
\Crefname{lemme}{Lemme}{Lemmes}
\crefname{definition}{définition}{définitions}
\Crefname{definition}{Définition}{Définitions}
\crefname{hypotheses}{hypothèses}{hypothèses}
\Crefname{hypotheses}{Hypothèses}{Hypothèses}
\crefname{exemple}{exemple}{exemples}
\Crefname{exemple}{Exemple}{Exemples}
\crefname{remarque}{remarque}{remarques}
\Crefname{remarque}{Remarque}{Remarques}
\crefname{note}{note}{notes}
\Crefname{note}{Note}{Notes}
\newcommand{\creflastconjunction}{ et~}
)
% ============================================================================

% ----------------------------------------------------------------------------
%  Polices et langue
% ----------------------------------------------------------------------------
\usepackage[T1]{fontenc}
\usepackage[utf8]{inputenc}     % à SUPPRIMER si lualatex/xelatex
\usepackage[french]{babel}
\usepackage{lmodern}
\usepackage{microtype}          % justification nettement meilleure

% ----------------------------------------------------------------------------
%  Géométrie
% ----------------------------------------------------------------------------
% 11 pt et 3 cm de marges latérales : environ 80 caractères par ligne,
% texte dense mais lisible.
\usepackage[a4paper,
            left=2.0cm, right=2.0cm, top=2.5cm, bottom=2.5cm,
            headheight=15pt]{geometry}

% ----------------------------------------------------------------------------
%  Mathématiques
% ----------------------------------------------------------------------------
\usepackage{amsmath, amssymb, amsthm, mathtools}
\usepackage{bm}
\usepackage{mathrsfs}

% ----------------------------------------------------------------------------
%  Flottants, tableaux, listes
% ----------------------------------------------------------------------------
\usepackage{graphicx}
\usepackage{float}
\usepackage{booktabs}
\usepackage{caption}
\usepackage{subcaption}         % APRÈS caption
\usepackage[inline]{enumitem}

% ----------------------------------------------------------------------------
%  Dessins
% ----------------------------------------------------------------------------
\usepackage{tikz}
\usetikzlibrary{arrows.meta, positioning, calc, angles, quotes,
    decorations.pathmorphing, patterns}

% ----------------------------------------------------------------------------
%  Couleurs, encadrés et liens
% ----------------------------------------------------------------------------
\usepackage{xcolor}
\definecolor{accent}{HTML}{1F3A5F}
\usepackage[most]{tcolorbox}

% Ligne « institution » de la page de titre.
\newcommand{\kicker}[1]{{\scshape\large\textcolor{accent}{#1}}}

\usepackage[colorlinks=true, linkcolor=accent, citecolor=accent,
            urlcolor=accent]{hyperref}
\usepackage{cleveref}           % IMPÉRATIVEMENT après hyperref

% ============================================================================
%  MISE EN PAGE
% ============================================================================

% --- 1. Titres ---------------------------------------------------------------
\usepackage{titlesec}

% Chapitre : « Chapitre » en petites capitales, gros numéro, titre
% en dessous, filet pour fermer.
\titleformat{\chapter}[display]
    {\normalfont}
    {%
        {\color{accent!55}\large\scshape\chaptertitlename}\;\;
        {\color{accent}\fontsize{52}{52}\selectfont\bfseries\thechapter}%
    }
    {10pt}
    {\color{accent}\Huge\bfseries}
    [\vspace{4mm}{\color{accent!40}\titlerule[1.4pt]}]
\titlespacing*{\chapter}{0pt}{20pt}{34pt}

% Structure : Chapitre -> I. -> 1. -> 1.1 -> a)
\setcounter{secnumdepth}{4}
\setcounter{tocdepth}{3}
\renewcommand{\thesection}{\Roman{section}.}
\renewcommand{\thesubsection}{\arabic{subsection}.}
\renewcommand{\thesubsubsection}{%
    \arabic{subsection}.\arabic{subsubsection}}
\renewcommand{\theparagraph}{\alph{paragraph})}

\titlespacing*{\section}{0pt}{2.2ex plus 1ex minus .2ex}{1.4ex plus .2ex}
\titlespacing*{\subsection}{0pt}{1.8ex plus 1ex minus .2ex}{1.0ex plus .2ex}
% a) : titre sur sa propre ligne
\titleformat{\paragraph}
    {\normalfont\normalsize\bfseries}{\theparagraph}{0.5em}{}
\titlespacing*{\paragraph}{0pt}{1.4ex plus .5ex minus .2ex}{0.6ex}

% --- 2. Flottants ------------------------------------------------------------
\renewcommand{\topfraction}{0.90}
\renewcommand{\bottomfraction}{0.75}
\renewcommand{\textfraction}{0.10}
\renewcommand{\floatpagefraction}{0.75}
\setcounter{topnumber}{3}
\setcounter{bottomnumber}{2}
\setcounter{totalnumber}{5}

% --- 3. Coupures de page -----------------------------------------------------
\widowpenalty=10000
\clubpenalty=10000
\raggedbottom

% --- 4. Listes et légendes ---------------------------------------------------
\setlist{itemsep=2pt, parsep=0pt, topsep=4pt}
\captionsetup{font=small, labelfont={bf,color=accent}, skip=6pt}

% --- 5. Accroche : phrase clé mise en exergue --------------------------------
% Usage : \begin{accroche} ... \end{accroche}
\newenvironment{accroche}[1][accent]{%
    \par\medskip\noindent
    \begingroup
    \color{#1}\itshape\large
    \rightskip=1.5em \leftskip=1.5em
    \parindent=0pt
    \hrule height 0.9pt\vspace{7pt}%
}{%
    \par\vspace{7pt}\hrule height 0.9pt
    \endgroup
    \par\medskip
}

% ============================================================================
%  NOTATIONS
% ============================================================================
% Vecteur : minuscule ou grecque, en gras italique
% (ex. \vect{u}, \vect{\omega}).
\newcommand{\vect}[1]{\bm{#1}}
% Matrice / tenseur : majuscule en gras droit (ex. \mat{R}, \mat{I}).
\newcommand{\mat}[1]{\mathbf{#1}}
% Quaternion : gras, souligné (ex. \quat{q} = (q_0, \vect{q})).
\newcommand{\quat}[1]{\underline{\bm{#1}}}
% Produit quaternionique : étoile centrée (\otimes est réservé au
% produit tensoriel).
\newcommand{\quatmul}{\ast}
% Repères : \rep{s} (sol), \rep{b} (solide).
\newcommand{\rep}[1]{\mathcal{R}_{#1}}
% Matrice antisymétrique associée à u : [u]_x v = u \wedge v
\newcommand{\hatx}[1]{[#1]_{\times}}
% Ensembles et opérateurs
\newcommand{\R}{\mathbb{R}}
\newcommand{\SO}{\mathrm{SO}(3)}
\newcommand{\Id}{\mat{I}}
\newcommand{\dd}{\mathrm{d}}
\newcommand{\ddt}{\frac{\dd}{\dd t}}
\newcommand{\pd}[2]{\frac{\partial #1}{\partial #2}}
\DeclareMathOperator{\rang}{rang}
\DeclareMathOperator{\Ker}{Ker}
\DeclareMathOperator{\Ima}{Im}

% ============================================================================
%  ENVIRONNEMENTS (numérotés par chapitre)
% ============================================================================
\theoremstyle{plain}
\newtheorem{theoreme}{Théorème}[chapter]
\newtheorem{corollaire}[theoreme]{Corollaire}
\newtheorem{proposition}[theoreme]{Proposition}
\newtheorem{lemme}[theoreme]{Lemme}
\theoremstyle{definition}
\newtheorem{definition}[theoreme]{Définition}
\newtheorem{hypotheses}[theoreme]{Hypothèses}
\newtheorem{exemple}[theoreme]{Exemple}
\theoremstyle{remark}
\newtheorem{remarque}[theoreme]{Remarque}
\newtheorem{note}[theoreme]{Note}
\renewcommand{\proofname}{Démonstration}
% Démonstration détaillée et preuve courte : même mise en forme
% (titre en italique, carré en fin).
\newenvironment{demonstration}[1][Démonstration]
    {\begin{proof}[#1]}{\end{proof}}
\newenvironment{preuve}[1][Preuve]
    {\begin{proof}[#1]}{\end{proof}}

% Couleurs des cadres : théorème en rouge, définition et hypothèses
% en vert, proposition, corollaire et lemme en bleu.
\definecolor{coulthm}{HTML}{A8322D}
\definecolor{couldef}{HTML}{2E7D32}
\definecolor{coulprop}{HTML}{1F5A9E}
\tcbset{
    cadre/.style={enhanced,colback=#1!7,colframe=#1!80,
        boxrule=0.6pt,leftrule=3pt,arc=1pt,left=6pt,right=6pt,
        before skip=6pt,after skip=6pt}}
\tcolorboxenvironment{theoreme}{cadre=coulthm}
\tcolorboxenvironment{definition}{cadre=couldef}
\tcolorboxenvironment{hypotheses}{cadre=couldef}
\tcolorboxenvironment{proposition}{cadre=coulprop}
\tcolorboxenvironment{corollaire}{cadre=coulprop}
\tcolorboxenvironment{lemme}{cadre=coulprop}
% Note : notation ou commentaire personnel, cadre gris.
\tcolorboxenvironment{note}{cadre=gray}

% Renvois \cref en français
\crefname{theoreme}{théorème}{théorèmes}
\Crefname{theoreme}{Théorème}{Théorèmes}
\crefname{corollaire}{corollaire}{corollaires}
\Crefname{corollaire}{Corollaire}{Corollaires}
\crefname{proposition}{proposition}{propositions}
\Crefname{proposition}{Proposition}{Propositions}
\crefname{lemme}{lemme}{lemmes}
\Crefname{lemme}{Lemme}{Lemmes}
\crefname{definition}{définition}{définitions}
\Crefname{definition}{Définition}{Définitions}
\crefname{hypotheses}{hypothèses}{hypothèses}
\Crefname{hypotheses}{Hypothèses}{Hypothèses}
\crefname{exemple}{exemple}{exemples}
\Crefname{exemple}{Exemple}{Exemples}
\crefname{remarque}{remarque}{remarques}
\Crefname{remarque}{Remarque}{Remarques}
\crefname{note}{note}{notes}
\Crefname{note}{Note}{Notes}
\newcommand{\creflastconjunction}{ et~}
